% ============================================================================
% CHAPTER 1: INTRODUCTION
% ============================================================================
\chapter{Introduction}
\label{ch:introduction}

\section{Background and Motivation}
\label{sec:background}

Recent climate regulations have led to a rise in the adoption of electric vehicles. More and more people are opting for an electric vehicle for their daily commute. According to studies, most of the time spent in commute is on the way to and from the office. Also, most people spend a third of their time in offices, and with them their cars, 
which can be charged during the workday. The car batteries can also be used in times of peak consumption to give electricity back to the building or to the grid. Since the charging will happen throughout the day,  
it is logical to create smart charging algorithms that can utilise information from solar power generation, pricing spreads throughout the day, and grid capacity limits in order to optimise the charging schedule. Office buildings also offer a huge rooftop space for the installation of solar panels; 
the energy produced by these panels can be used both for the building's consumption as well as for the charging of the electric vehicles. The cost of electricity produced is zero, and this helps reduce the overall cost. As such, the utilization of renewable energy is the most important aspect of the charging algorithm. 

\section{Literature Review}
\label{sec:literature_review}

Currently, there are many different algorithms that are used for the smart scheduling of the EV fleet. Early studies, such as Richardson \cite{richardson2012ev_ucd}, used linear programming in order to decide which vehicle will be charged and when in 15-minute intervals,
taking into account network limitations. Further studies included the use of reinforcement learning to increase grid stability, and optimize for pricing and renewables utilization; several studies are included in this review \cite{qiu2023rl_ev_review}.

\section{Smart Charging Levels}
\label{sec:smart_charging_levels}

The algorithms developed in this thesis are separated by their intelligence levels. Level 0 is a simple rule-based algorithm that prioritises solar renewable energy, respecting grid constraints and the target charging SOC of the EV battery; it charges without respecting future price levels. 
Level 1 is a Q-learning algorithm that uses a reward function to include information about price levels and also the availability of solar energy. Level 2 is a Q-learning algorithm that also includes the action of discharging energy back to the grid. The remaining chapters focus on different parts of the algorithms.

\section{Thesis Structure}
\label{sec:structure}

Chapter~2 focuses more on the models that describe the office, the EVs, the grid, and the solar panels. Chapter~3 focuses on the data for the office building consumption, electricity pricing, and solar generation. Chapter~4 focuses on the simulation using each algorithm, collecting and comparing the results.
Chapter~5 focuses on the conclusion of the results and recommendations for future work.